\chapter{About the Dataset}

The dataset used in the Kaggle competition \parencite{raoDepressionSurveyDataset} includes a range of predictors relevant to mental health. For the purposes of our fairness audit, we consider \texttt{Gender} as the primary sensitive attribute. In addition to gender, features such as profession, degree, and dietary habits may act as proxy variables, potentially capturing implicit gender-related signals.

Beyond individual attributes, the model’s performance may also be influenced by geographic and cultural biases embedded in the dataset. The data originates from the Indian subcontinent, but distribution plots
[\href{https://colab.research.google.com/drive/1V5NlYPUzOfSxUoDHg67yVfatMXSTWNSU?usp=sharing}{1}, \href{https://www.google.com/maps/d/edit?mid=17YtHn2h6j22KFGfMiAvcx5rToAd4Ukw&usp=sharing}{2}] suggest it is disproportionately sampled from major urban centers in North and Central-West India. Areas such as smaller cities, rural regions, and the South, Central-East, and North-East are notably underrepresented. This sampling imbalance raises concerns about whether the dataset reflects the full regional diversity of India, and whether the trained model can generalize well across culturally and geographically distinct populations.

Importantly, the dataset is \textbf{synthetically generated}. As a result, biases in regional distribution may stem from (i) the original data used for synthesis, (ii) the design of the generation algorithm, or (iii) deliberate sampling decisions. Yet, there appears to be no clear rationale for the selection of included cities. The sampled regions do not consistently correspond to city size, population density, state capitals, metropolitan zones, or any coherent regional grouping. This absence of justification limits confidence in the model’s representativeness and generalizability.

Furthermore, the dataset omits important cultural and familial dimensions that are particularly salient in the Indian context. In many Asian societies, family structure, interdependence, and relationship quality play central roles in shaping psychological well-being. For instance, Lu and colleagues \parencite{lu2010understanding} demonstrate that even the conceptualization of depression differs significantly between Chinese and American populations, reinforcing the importance of culture in mental health modeling.

To that end, variables capturing intergenerational living, number of dependents, or relationship satisfaction would have been valuable additions to the dataset. Their inclusion could offer more culturally grounded insights into depression risk, especially for a demographically diverse population like India’s. Without these features, the model may overlook key factors that contribute to mental health outcomes in the region.

See Appendix~\ref{appendix} for further details on data distribution and regional representation.